% ! TeX root = ../main-ieee.tex

\section{Introduction}\label{sec:amr:intro}

Block-structured Adaptive Mesh Refinement (AMR) codes are critical to modern scientific computing,
enabling large-scale simulations in astrophysics, cosmology, climate science, and fluid
dynamics~\cite{Dubey2014:Survey, Burst2011:p4est}.
By dynamically refining the mesh around steep gradients, AMR delivers 1--2 orders of magnitude
performance gains over uniform approaches~\cite{Vitti:DecadeAMR}.
These codes often run for weeks on hundreds of thousands of
cores~\cite{Liska2022:H-AMR,Burst2011:p4est,Grete2023:Parthenon}, making performance optimization
essential.
Load balancing remains a major challenge~\cite{Dubey2014:Survey,Burst2008:ALPS}: while octree and
space-filling curve (SFC) placements work for uniform workloads~\cite{Burst2011:p4est}, they do not
account for compute kernel variability.
% Stragglers in parallel codes are costly due to frequent global synchronization, and the
% fine-grained adaptivity of AMR exacerbates this effect.
Parallel codes are limited by stragglers---ranks lagging significantly behind peers---whose
delay propagates due to frequent global synchronization, and the fine-grained adaptivity of AMR
exacerbates this effect.
Prior analyses of proxy application traces~\cite{Klenk2017:OverviewMPICharacteristics} document
significant MPI waiting time (\textgreater{}60\% at 1000 ranks), while noting the limitations of
standard trace data for the deeper, cross-stack diagnosis required to pinpoint root causes.
The root causes, whether algorithmic or infrastructural, remain entangled.
We ask: what changes when performance instrumentation has full visibility over a real AMR code?

To answer this, we sought to improve AMR placement using fine-grained runtime telemetry on a
research cluster providing full-stack access.
Our goal was to replace static heuristics with telemetry-driven workload models.
We adopted an integrated workflow that combined telemetry collection, analysis, and intervention.
However, initial experiments immediately revealed a deeper challenge: observed metrics often
diverged unexpectedly from predictors such as message volume, undermining our ability to construct
reliable workload models.
Rather than directly enabling better placement, fine-grained telemetry first required us to
diagnose and mitigate runtime artifacts to make the data useful for placement decisions
(\S\ref{sec:amr:chal}).
% to create full access required us to first interpret and diagnose our stack for predictable te
% first exposed the fundamental difficulty of interpreting cross-layer performance interactions
% (\S\ref{sec:amr:chal}).

This paper describes the end-to-end process we undertook to develop telemetry-driven placement
policies.
We began by systematically diagnosing variability across the stack and tuning system parameters.
This required extensive analysis of large volumes of telemetry, prompting a shift from standard
tracing and profiling tools to a custom analytical pipeline over structured data
(\S\ref{sec:amr:tele}).
Tuning the system stack resulted in predictable runtime response as placement aspects such as load
balance and locality were varied.
Recognizing that these two properties represented competing optimization objectives, we developed
CPLX (\S\ref{sec:amr:design}), a hybrid placement policy balancing these two objectives via a tunable
parameter to minimize overall runtime.
Because AMR codes refine frequently, placement must meet strict timing constraints ($<$ 50ms in our
target codes).
CPLX meets these requirements while leveraging existing octree and SFC infrastructure for
compatibility with current AMR frameworks.

CPLX improves runtime by up to 21.6\% over optimized baselines (\S\ref{sec:amr:eval}) through its
tunable control over the load-locality tradeoff, validating our process-driven methodology.
However, our experience suggests that performance optimization in large-scale systems is inherently
context-dependent.
While algorithmic approaches are necessary, their application must be embedded in a broader
empirical process---one that interprets telemetry in the context of the runtime stack and prioritizes
objectives only insofar as they measurably impact end-to-end performance.
Accordingly, we focus not just on CPLX itself, but on the end-to-end process that led to it, and
distill methodological lessons applicable to other large-scale tuning efforts.
This paper makes three key contributions:

\begin{itemize}
  \item An empirical case study documenting the challenges of achieving reliable,
        interpretable telemetry for AMR codes using full-stack access and iterative
        tuning.
  \item A novel placement policy, CPLX, that provides tunable control over the tradeoff
        between compute balance and communication locality (\S\ref{sec:amr:design}).
  \item Broader lessons on telemetry-driven performance analysis and the empirical,
        context-specific nature of tuning in large-scale parallel systems.
\end{itemize}

We provide background on block-structured AMR codes and current placement approaches
(\S\ref{sec:amr:bg}), then outline the challenges of incorporating telemetry into placement decisions
(\S\ref{sec:amr:chal}).
We present our methodology for obtaining reliable performance measurements (\S\ref{sec:amr:tele}),
followed by the design of our placement policies (\S\ref{sec:amr:design}) and their evaluation
(\S\ref{sec:amr:eval}).
We then summarize the broader lessons drawn from our experiences (\S\ref{sec:amr:lessons}) before
concluding with related work (\S\ref{sec:amr:relwork}).

All our policies and associated artifacts are open source~\cite{Ankus2025:cplx}.